% HD-S42. GPT-6 Astra Pro /root/sixfolds; inherited effort not exposed.
% 21 November 2026. Actual support and field descent, all ranks.
\subsection{HD-S42: supports of a reducible normal bireflection subgroup}
\label{sec:reducible-bireflection-support}

Let \(K\subset\mathbb C\) have characteristic zero,
let \(G\subset GL_K(V)\) be finite and absolutely
irreducible, and let \(1\ne N\triangleleft G\) be
generated by actual rank-two elements. Take all
\(G\)-conjugates of a generating set, and call this
\(G\)-stable set \(\mathcal X\). The following is an
exact structural reduction over the specified field \(K\).

Either the complex \(N\)-isotypic components all have
multiplicity one, or they all have multiplicity two. In the
first case every isotypic component is \(K\)-defined, or
the components pair into \(K\)-defined blocks and are
exchanged by one quadratic extension of \(K\). In the
second case every isotypic component is \(K\)-defined
and has a literal tensor form \(W_K\otimes_K K^2\),
where the corresponding subgroup of \(N\) is an actual
irreducible \(K\)-reflection group on \(W_K\).
All these systems are permuted transitively by \(G\).
The maps and the paired case's exact quadratic field are
proved below.

\paragraph{1. Isotypic projectors and the multiplicity bound.}
Average any projection over the finite group \(N\) to obtain
invariant complements, and write the complete complex
isotypic decomposition as
\[
 V_{\mathbb C}=\bigoplus_{i\in I}U_i,
       \qquad U_i=W_i\otimes\mathbb C^e,
       \qquad \dim W_i=r.                                  \tag{S42.1}
\]
Here \(N\) acts on the first tensor factor. Normality
makes \(G\) permute its irreducible types. Any proper
union of orbits would give a proper \(G\)-invariant sum,
so the action on \(I\) is transitive. The dimensions
and multiplicities are consequently common \(r,e\),
as displayed. This proves the needed normal-subgroup
decomposition directly.

For the character \(\chi_i\) on \(W_i\), the actual
isotypic projector is
\[
      p_i=\frac r{|N|}\sum_{n\in N}
                       \chi_i(n^{-1})\rho(n).              \tag{S42.2}
\]
Character averaging, including the multiplicity factor,
proves that these are exactly the complete primitive central
idempotents of \(\mathbb C\rho(N)\). They lie in
\(L\rho(N)\), for a finite cyclotomic Galois extension
\(L/K\). Every Galois automorphism permutes the \(p_i\)
and commutes with the permutation of \(G\), since all
the original \(G\)-matrices have entries in \(K\).

For \(x\in\mathcal X\), put
\(S_x=\{i:(x-1)p_i\ne0\}\). Rank in (S42.1) gives
\[
              2=e\sum_{i\in I}
                   \operatorname{rank}(\rho_i(x)-1).       \tag{S42.3}
\]
Thus \(e=1\) or \(2\). Every \(S_x\) is individually
fixed by Galois: apply it to the literal matrix equation
\((x-1)p_i=0\). The collection of these supports is
permuted by \(G\). Every label occurs in a support;
otherwise transitivity would make \(N\) act trivially
everywhere, contrary to \(N\ne1\).

\paragraph{2. Multiplicity one: single types or matched pairs.}
Let \(e=1\). The supports have size one or two. If a
singleton occurs, its \(G\)-conjugates give every singleton,
and Galois fixes every label. If none occurs, regard the
supports as the edges of a graph on \(I\). This graph
has no isolated vertex and is vertex-transitive under
\(G\). If its degree is at least two, each vertex is
the intersection of two individually fixed edges. Again
Galois fixes every vertex. In both situations
\[
                   p_i\in M_m(K),\qquad U_i=p_iV_K
                                                               \tag{S42.4}
\]
are actual \(K\)-subspaces, permuted transitively by \(G\).

The only remaining graph is a perfect matching. Each pair
sum \(p_i+p_{i'}\) is Galois fixed. The Galois image on
all labels has order at most two. Indeed a permutation
commuting with transitive \(G\) and fixing one label fixes
every label. A nontrivial Galois permutation must therefore
have no fixed label; since it fixes every matching edge, it
must swap the endpoints of every edge simultaneously.
If this image is trivial, (S42.4) still applies. Otherwise
let \(E\) be the fixed field of its kernel in \(L\).
Then \([E:K]=2\), every \(p_i\) lies in \(M_m(E)\),
and the nonidentity \(\tau\in\operatorname{Gal}(E/K)\)
exchanges precisely each matched pair.

For one pair put \(U=(p+p')V_K\). The explicit field
action is
\[
 \iota:E\longrightarrow\operatorname{End}_K(U),\qquad
                     a\longmapsto ap+\tau(a)p'.             \tag{S42.5}
\]
The displayed matrix is fixed by Galois and hence has entries
in \(K\); it is multiplicative, sends one to \(I_U\),
and is injective. Moreover
\[
 U\longrightarrow p(V\otimes_K E),\quad u\longmapsto pu,
 \qquad w\longmapsto w+\tau(w)                              \tag{S42.6}
\]
are inverse \(E\)-linear maps for (S42.5). They make \(U\)
an actual \(E\)-space of dimension \(r\).

Every generator supported on this edge has rank one on
each of its two complex types. Generators on distinct matching
edges commute and have disjoint supports, so \(N\) is the
direct product of its actual supported subgroups. On \(U\)
the corresponding image is faithful and \(E\)-linear,
generated by \(E\)-reflections. Faithfulness on the first
\(E\)-component follows from conjugation: identity on that
component implies identity on its conjugate, hence on \(U\).
Its character \(\chi\) has values in \(E\). If all
those values lay in \(K\), the two conjugate characters
would be equal, contrary to the inequivalence of the two
types. For its actual character field \(F\), therefore,
\[
                               KF=E.                       \tag{S42.7}
\]
The pair stabilizer in \(G\) acts semilinearly on (S42.5),
normalizes this reflection-group image, and includes an
embedding flip because \(G\) is transitive on individual
types. Equations (S42.5)--(S42.7) are the precise paired
coefficient data used in HD-G30.

\paragraph{3. Multiplicity two: actual tensor coordinates.}
Let \(e=2\). Equation (S42.3) says that every generator
has support exactly one type, and acts as a reflection on
its \(W_i\). For a generator \(x\) supported at \(i\),
the literal subspace
\[
        U_i(K)=K[N]\operatorname{im}_K(x-1)\subset V_K
                                                               \tag{S42.8}
\]
has complexification exactly \(W_i\otimes\mathbb C^2\).
Indeed \(\operatorname{im}(x-1)\) is a nonzero line in
\(W_i\) tensor the full multiplicity space; its \(N\)-span
is the entire displayed tensor product by irreducibility.
Thus these components are individually \(K\)-defined.
The supported subgroups again form a literal direct product.

On one such component \(U\), let \(A\) be the span over
\(K\) of its supported subgroup. Its complexification is
\(M_r(\mathbb C)\otimes I_2\), by matrix-coefficient
averaging. It is a central simple \(K\)-algebra. There is
no division-algebra obstruction here, as the following maps
show explicitly. Since \(U\) is \(K\)-defined, the
irreducible character equals one half of its trace and has
values in \(K\). For a supported generator \(x\), its
nonidentity eigenvalue on \(W_i\) is
\(\lambda=\chi_i(x)-r+1\in K\). Consequently
\[
          P=\frac{x-I_U}{\lambda-1}\in A,\qquad
          P^2=P,\qquad\dim_K PU=2,\qquad PAP=KP.           \tag{S42.9}
\]
The last identity follows either by its scalar action on the
complex reflection line or by scalar extension of the corner.
Put \(W_K=AP\) and \(M_K=PU\). Left multiplication
identifies \(A\) with \(\operatorname{End}_K(W_K)\):
its kernel is an ideal annihilating \(AP\), hence zero
since \(APA=A\), and the explicit matrix-unit description
after extension has \(\dim_K W_K=r\), \(\dim_K A=r^2\).
The actual tensor map is
\[
 W_K\otimes_K M_K\longrightarrow U,
                    aP\otimes u\longmapsto au.              \tag{S42.10}
\]
Choose \(a_j,b_j\in A\) with \(\sum_j a_jPb_j=1\),
which exist because \(APA=A\). Its inverse is
\(u\mapsto\sum_j a_jP\otimes Pb_j u\); the identity
\(PAP=KP\) verifies both composites. The supported
subgroup acts faithfully by left multiplication on \(W_K\),
as an absolutely irreducible \(K\)-reflection group, and
trivially on \(M_K\).

Let \(H\) stabilize \(U\). Its representation on \(U\)
is absolutely irreducible: a proper \(H\)-invariant subspace
would induce, by the actual component translates, a proper
\(G\)-invariant subspace. It normalizes \(A\). Every
automorphism of the explicit matrix algebra \(A\) is inner
over \(K\), with a concrete construction: for its transformed
matrix units choose \(u_1\ne0\) in the image of the first
diagonal unit and take \(u_j\) to be the image of \(u_1\)
under the \((j,1)\) unit. These columns form the required
intertwining matrix. After removing that matrix on \(W_K\),
the remaining action centralizes \(A\) and acts on \(M_K\).
Thus (S42.10) gives actual homomorphisms
\[
                 H\longrightarrow PGL(W_K),\qquad
                 H\longrightarrow PGL(M_K),\quad\dim M_K=2.
                                                               \tag{S42.11}
\]
No tensor decomposition has been inferred merely from an
equality of dimensions.

\paragraph{4. Cover and primitive-factor maps.}
For this paragraph assume the imaginary-quadratic sixfold
setting, with \(K\) the exact character field of the
original coefficient representation. The structural assertions
of paragraphs1--3 retain their stated general-field scope.
For any descended block \(U\) above, let \(H\) be its
stabilizer, \(\sigma:H\to GL_K(U)\) its actual action,
and \(H_0=\ker\sigma\). The block-translates map
\(\operatorname{Ind}_H^G\sigma\to V\),
\(t_j\otimes u\mapsto\rho(t_j)u\), is an isomorphism.
For a connected original cover put
\[
       Y=\widetilde C/H,\qquad C'=\widetilde C/H_0,
                         \alpha:\widetilde C\to C'.        \tag{S42.12}
\]
The exact Shapiro functions are
\(F_v(x)=\pi_U\rho(x)v\), with
\(F(hx)=\sigma(h)F(x)\), right translation by \(G\),
and inverse \(F\mapsto\sum_{H\backslash G}\rho(x)^{-1}F(x)\).
The label \(Ht_j^{-1}\) is the block \(t_jH\).
These are the finite-pushforward and invariant-stalk maps of
HD-S30 and HD-S36. The coefficient forms make these components
orthogonal: their irreducible \(N\)-types are disjoint.
The rational matrix-coefficient corner idempotents for a vector
in \(U\) satisfy \(e_G=e_H\), since the coefficients
off \(H\) vanish and
\(\dim V/|G|=\dim U/|H|\). Thus pullback gives the actual
primitive-factor isogeny, with
\[
       \alpha_*\alpha^*=[|H_0|],\qquad
       (\alpha^*)^*\theta_{J\widetilde C}=|H_0|\theta_{JC'}.
                                                               \tag{S42.13}
\]
The prescribed original corner is carried to this one by the
literal \(b_{v,w}\) and integer-cleared isogeny in
(S41.15)--(S41.16). These maps retain the corner field,
period rank, and ordinary rank-six Hermitian determinant,
since each positive rational multiplier has sixth power a
quadratic-field norm. The induced trace formula also proves
that a descended block has the same exact character field
\(K\) as the original representation.

\paragraph{5. Consequences for a rank-nine sixfold family.}
Assume now the same imaginary-quadratic sixfold conventions
as HD-S41. The determinant applications to HD-S36, HD-G27,
and HD-S34 in this paragraph are restricted to
\[
 K=\mathbb Q(i)\quad\hbox{or}\quad
 K=\mathbb Q(\sqrt{-d}),\qquad
 d\in\{2,7,11,19,43,67,163\}.
\]
The coefficient, genus, and tensor bounds themselves retain
their scope over every imaginary quadratic \(K\).
A descended coefficient block of rank \(q\ge2\)
has
\[
                6=q(2g(Y)-2)+\sum c'_j.                   \tag{S42.14}
\]
For \(g(Y)=1\), the marked-base dimension is at most six.
For \(g(Y)=2\), one has \(q\le3\) and dimension at
most \(3+6-2q\le5\). For \(g(Y)\ge3\) the right
side is at least eight. Thus a rank-nine period family with
\(q\ge2\) has \(Y\simeq\mathbb P^1\), at least
twelve true coefficient branch values, and total codimension
\(2q+6\).

In the individually descended multiplicity-one case, the
supported image of \(N\) on a block is absolutely
irreducible and is generated by elements of rank one or two.
Any rank-one element invokes exactly HD-S36 or HD-G27; if
none occurs and \(q>8\), HD-S41 applies to that actual
normal subgroup of the faithful block group. For \(2\le q\le8\)
the branch budget itself forces a reflection. Rank two is
also excluded directly by the period bound seven. A rank-one
block is the actual induced character of HD-S34 in the
Gaussian field, while roots of unity \(\{\pm1\}\) would
give rational induced trace. Hence this descended case yields
only the previously covered determinant classes, with the
literal transports (S42.12)--(S42.13).

In the multiplicity-two case, write \(q=2r\). Rank one
for \(W_K\) would give \(q=2\), already excluded. For
\(r\ge2\), the projective image on \(M_K\) in
(S42.11) is noncyclic: a cyclic finite projective image has
a diagonalizable generator and a common eigenline, which
would make \(W_K\) tensor that line invariant under \(H\).
For the generating sphere inertia tuple there must be at
least three nonidentity images in this projective factor;
zero or one is impossible, and two would be inverse and
generate a cyclic image. Each such tensor element has fixed
dimension at most \(r\), hence codimension at least \(r\).
Every other nonidentity tensor has codimension at least two.
For \(b\ge12\), therefore,
\[
                4r+6\ge3r+2(b-3)\ge3r+18,
                              \qquad r\ge12.               \tag{S42.15}
\]
This is the exact surviving tensor range, not an assertion
that any tuple or period-rank-nine source in that range exists.

The matched-pair endgame uses the explicit quadratic reflection
data (S42.5)--(S42.7) and is treated separately in HD-G30.
The multiplicity-two endgame starts from the actual reflection
tensor factor (S42.9)--(S42.11). The support theorem and all
the consequences stated here are proved independently of any
unproved conclusion about those remaining endgames.
